% Uninput; independent literal audit requested.
\subsection{Two excess arrows cannot support second cohomology}

\begin{lemma}\label{C:arrow-excess-two-H2}
Let $C$ be a connected finite category with $k$ objects and
$n\leq2k+1$ arrows.  Then $H^2(BC;\mathbb Q)=0$.
\end{lemma}
\begin{proof}
Put $e=n-(2k-1)\leq2$.  Deleting an object which is a
retract of another preserves the classifying-space
homotopy type and connectedness, removes at least its
identity and two cross arrows, and lowers $e$ by at
least one.  Repeat until no such object remains.
If one object remains, its monoid has at most three
elements and has zero second rational cohomology by
Theorem~\ref{E:all-six-types}.

A nontrivial strongly connected component with $t$
objects has at least $t^2+t$ arrows: all cross hom-sets
are nonempty, and every corner has at least two
elements, since a trivial corner would give a retract.
If it is the whole category, its excess is at least
$t^2-t+1\geq3$.  If it is proper, the connected support
budget in the proof of
Theorem~\ref{C:arrow-excess-four} gives excess at least
$t^2\geq4$.  Both are impossible, so support is acyclic.

Write $D,P,R$ for the corner, parallel-arrow and support
cycle excesses as in
\eqref{C:excess-decomposition}.  Then
$D+P+R=e\leq2$.
If $P=0$, the actual support-poset functor induces an
isomorphism on $H^2(-;\mathbb Q)$ by
Lemma~\ref{C:singleton-support-H2}.  Its comparability
graph has cycle rank at most two.  It can contain at
most two triangles: the boundaries of any three
distinct triangles are linearly independent, because
each triangle has an edge absent from the other two.
The boundaries of one or two triangles are independent
as well.  Thus the degree-two simplicial boundary is
injective, and the support poset has zero second
homology.

Suppose $P\geq1$.  A nonsemisimple corner algebra would
require a monoid of order at least three and hence
$D\geq2$, contradicting the budget.  All corner
algebras are therefore semisimple.  If $R=0$, the
support tree has no composable strict pair, so the
relative bar has length at most one and has zero
second homology.  Otherwise equality is forced:
$D=0$, $P=1$, and $R=1$.  All corners are trivial and
there is only one doubled hom-set.  A unicyclic support
without a triangle has no two-simplices.  If its unique
cycle is a triangle, the nerve is the union of that
triangle category and its graph bridges, meeting at
points.  At least one adjacent hom-set in the triangle
is a singleton.  The component of its degree-two
boundary on the other adjacent edge is injective on
the two-simplex labels.  Thus its second homology is
zero, as is that of every bridge.  The union consequently
has zero second rational homology.
Duality completes the proof.
\end{proof}
